%! Introducing Statistics: What Can We Learn from Data? — Unit 1, Lesson 1.0 — Student Packet Cover Sheet
\documentclass[10pt]{article}
\usepackage{apstats-article}
\usepackage{apstats-boxes}
\usepackage{ltablex}
\keepXColumns

\begin{document}

% Full-bleed navy banner
\begin{tikzpicture}[remember picture, overlay]
  \fill[navy] (current page.north west)
    rectangle ([yshift=-0.9in]current page.north east);
\end{tikzpicture}
\vspace{-0.8in}

\begin{center}
  {\color{white}\textbf{\LARGE Statistics}}\\[4pt]
  {\color{skymid}\large\textbf{Unit 1}}\\[6pt]
  {\color{white}\textbf{\large Lesson 1.0 \quad Introducing Statistics: What Can We Learn from Data?}}
\end{center}

\vspace{0.22in}
\namedateperiod
\vspace{0.18in}

\boxguard[14]
\begin{learningtargetbox}
\small
% SPOILER RULE (EFFL): plain language only. The formal terms — data, individual,
% variable, value, variability, statistical question — are what the debrief attaches,
% so none of them may appear here. The AP Learning Objective (VAR-1.A) lives in the
% teacher-facing lesson plan.
\begin{enumerate}[leftmargin=1.5em, itemsep=5pt]
  \item \ldots{} explain why a number on its own tells you nothing, and say what you would
        need to be told before it does.
  \item \ldots{} read a table of recorded information and say what one row stands for and
        what each column keeps track of.
  \item \ldots{} tell apart the questions you can answer by looking one thing up and the
        questions that need information about a whole group.
  \item \ldots{} write a question of my own that a set of recorded information could answer.
\end{enumerate}
\end{learningtargetbox}

\vspace{0.14in}

\boxguard[14]
\begin{tocbox}
\small
\renewcommand{\arraystretch}{1.5}
\begin{tabularx}{\linewidth}{@{} c l X r @{}}
\rowcolor{sky}
\textbf{\#} & \textbf{Component} & \textbf{Description} & \textbf{Score} \\
\midrule
1 & Warm-Up & Seven numbers with no context --- what can you tell? & \blank{1.2cm} \\
2 & Experience \& Formalize & \emph{The Adoption Board} --- two problems with your group, then
                 the QuickNotes we fill in together and one application we work together
                 & \blank{1.2cm} \\
% Check Your Understanding is PRACTICE and is NOT scored: the score cell prints NA, never a
% \blank{}. The graded practice in this course is the homework row below it.
3 & Check Your Understanding & Five exercises in new contexts ---
                 \emph{these are for practice, so they are not scored} & \textbf{NA} \\
4 & Homework & The adoption board again: AP-style multiple choice and one free-response set
                 & \blank{1.2cm} \\
\midrule
  & \multicolumn{2}{r}{\textbf{Total}} & \blank{1.2cm} \\
\end{tabularx}
\end{tocbox}

\vspace{0.14in}

\boxguard[8]
\begin{remindbox}
\small
In this lesson you will \textbf{experience first, formalize later}: you and your group will work
out the adoption board using only what you already know --- and only afterward will we name what
you found. Do not wait to be told the words. Write what you see, and be ready to point at the
board and say \emph{how} you know.
\end{remindbox}

\end{document}
